% ------------------------------------------------------------------
% Sección de Washington Apunte Pazmiño
% Bloque: publicación con imagen, distribución al muro, seguidores
% y notificaciones (requisitos R3 y R4; situación de fallo F2).
% Integración en notas.tex: % ------------------------------------------------------------------
% Sección de Washington Apunte Pazmiño
% Bloque: publicación con imagen, distribución al muro, seguidores
% y notificaciones (requisitos R3 y R4; situación de fallo F2).
% Integración en notas.tex: % ------------------------------------------------------------------
% Sección de Washington Apunte Pazmiño
% Bloque: publicación con imagen, distribución al muro, seguidores
% y notificaciones (requisitos R3 y R4; situación de fallo F2).
% Integración en notas.tex: % ------------------------------------------------------------------
% Sección de Washington Apunte Pazmiño
% Bloque: publicación con imagen, distribución al muro, seguidores
% y notificaciones (requisitos R3 y R4; situación de fallo F2).
% Integración en notas.tex: \input{secciones/publicacion-muro}
% ------------------------------------------------------------------

\subsection{Publicación y distribución al muro}

\paragraph{Recorrido de un aviso.}
El usuario o el negocio local crea el aviso desde el cliente (CLI). La primera llamada
(S1) envía a PUB el \texttt{id\_usuario}, el texto, la categoría, la ubicación y el tipo
de aviso (normal o alerta); PUB lo registra y responde con el \texttt{id\_aviso}. Si el
aviso lleva imagen, el cliente la sube directamente al servicio externo de almacenamiento
(S2), que devuelve la URL o el identificador del archivo, y después registra esa
referencia en PUB (S2b). Así el archivo no pasa por los servicios internos.

Cuando el aviso está completo, PUB publica el evento \texttt{AvisoPublicado} en Kafka (A1)
y no espera nada más: el usuario ya recibió su confirmación. MUR consume el evento (A2),
consulta a SEG los seguidores del autor mediante gRPC (S3) y agrega el aviso al muro de
cada seguidor. En paralelo, NOT consume el evento \texttt{NuevaPublicación} o
\texttt{AlertaSeguridad} (A3) y envía la notificación push a los seguidores (A5). Cuando un
seguidor abre la aplicación, obtiene de MUR su muro ya armado (S10).

\begin{center}
\small
\setlength{\tabcolsep}{4pt}
\begin{tabular}{|l|l|l|l|l|}
\hline
\textbf{ID} & \textbf{Origen $\rightarrow$ destino} & \textbf{Tipo} & \textbf{Sincronía} & \textbf{Protocolo y middleware} \\
\hline
S1  & CLI $\rightarrow$ PUB             & Petición-respuesta & Síncrona  & HTTPS/TCP, REST API \\
S2  & CLI $\rightarrow$ almac. externo  & Petición-respuesta & Síncrona  & HTTPS/TCP, API externa \\
S2b & CLI $\rightarrow$ PUB             & Petición-respuesta & Síncrona  & HTTPS/TCP, REST API \\
A1  & PUB $\rightarrow$ Kafka           & Mensajería         & Asíncrona & TCP, Kafka \\
A2  & Kafka $\rightarrow$ MUR           & Mensajería         & Asíncrona & TCP, Kafka \\
S3  & MUR $\rightarrow$ SEG             & RPC                & Síncrona  & HTTP/2, gRPC \\
A3  & Kafka $\rightarrow$ NOT           & Mensajería         & Asíncrona & TCP, Kafka \\
A5  & NOT $\rightarrow$ CLI             & Mensajería push    & Asíncrona & HTTPS/TCP, FCM \\
S10 & CLI $\rightarrow$ MUR             & Petición-respuesta & Síncrona  & HTTPS/TCP, REST API \\
\hline
\end{tabular}
\end{center}

\paragraph{Mecanismos y justificación.}
\begin{itemize}[leftmargin=*]
  \item \textbf{REST sobre HTTPS, síncrono (S1, S2, S2b).} Crear el aviso, subir la imagen
  y vincularla son acciones en las que el usuario necesita una respuesta inmediata
  (\texttt{id\_aviso} o URL) para continuar. HTTPS funciona de forma nativa en Android y en
  el navegador y protege los datos en tránsito.
  \item \textbf{Kafka, asíncrono (A1, A2, A3).} La distribución al muro y las
  notificaciones pueden completarse después de responder al usuario. Kafka desacopla a PUB
  de MUR y de NOT: PUB publica el evento y continúa, y el evento queda guardado aunque un
  consumidor esté caído.
  \item \textbf{gRPC sobre HTTP/2, síncrono (S3).} MUR necesita cada página de seguidores
  para seguir distribuyendo. gRPC trabaja con mensajes binarios compactos y un contrato
  \texttt{.proto} entre servicios internos, por lo que tiene menos latencia que REST con
  JSON.
\end{itemize}

\paragraph{R3: tiempos distintos para avisos normales y alertas.}
Un aviso normal puede tardar hasta 5~s en aparecer en el muro de los seguidores, así que se
acepta consistencia eventual. Una alerta de seguridad debe aparecer en menos de 1~s. Para
que una alerta no quede en cola detrás de los avisos normales, PUB la publica en un tópico
propio (\texttt{alertas.seguridad}) con consumidores dedicados en MUR y en NOT. La alerta
tiene prioridad y capacidad de procesamiento separada del resto.

\paragraph{R4: un negocio con 30.000 seguidores.}
PUB no espera a que el aviso llegue a los 30.000 seguidores, porque la distribución es
asíncrona (A1). MUR pide los seguidores a SEG por páginas (S3) y reparte el trabajo en
lotes entre varios consumidores en paralelo. Además, se limita la carga que puede generar
un solo autor, para que una publicación masiva no ocupe toda la capacidad y el servicio
siga respondiendo al resto de usuarios.

\paragraph{F2: el servicio de muro se cae mientras un negocio publica.}
La distribución no se pierde. MUR confirma un evento en Kafka (\emph{commit} del
\emph{offset}) solo después de procesarlo; si se cae a mitad del trabajo, los eventos no
confirmados siguen en Kafka y, al recuperarse, MUR los vuelve a leer desde el último punto
confirmado. Como un evento puede procesarse dos veces, el consumidor es idempotente:
agregar el mismo aviso al mismo muro no crea un duplicado, por ejemplo gracias a una clave
única (\texttt{id\_usuario}, \texttt{id\_aviso}).

% ------------------------------------------------------------------

\subsection{Publicación y distribución al muro}

\paragraph{Recorrido de un aviso.}
El usuario o el negocio local crea el aviso desde el cliente (CLI). La primera llamada
(S1) envía a PUB el \texttt{id\_usuario}, el texto, la categoría, la ubicación y el tipo
de aviso (normal o alerta); PUB lo registra y responde con el \texttt{id\_aviso}. Si el
aviso lleva imagen, el cliente la sube directamente al servicio externo de almacenamiento
(S2), que devuelve la URL o el identificador del archivo, y después registra esa
referencia en PUB (S2b). Así el archivo no pasa por los servicios internos.

Cuando el aviso está completo, PUB publica el evento \texttt{AvisoPublicado} en Kafka (A1)
y no espera nada más: el usuario ya recibió su confirmación. MUR consume el evento (A2),
consulta a SEG los seguidores del autor mediante gRPC (S3) y agrega el aviso al muro de
cada seguidor. En paralelo, NOT consume el evento \texttt{NuevaPublicación} o
\texttt{AlertaSeguridad} (A3) y envía la notificación push a los seguidores (A5). Cuando un
seguidor abre la aplicación, obtiene de MUR su muro ya armado (S10).

\begin{center}
\small
\setlength{\tabcolsep}{4pt}
\begin{tabular}{|l|l|l|l|l|}
\hline
\textbf{ID} & \textbf{Origen $\rightarrow$ destino} & \textbf{Tipo} & \textbf{Sincronía} & \textbf{Protocolo y middleware} \\
\hline
S1  & CLI $\rightarrow$ PUB             & Petición-respuesta & Síncrona  & HTTPS/TCP, REST API \\
S2  & CLI $\rightarrow$ almac. externo  & Petición-respuesta & Síncrona  & HTTPS/TCP, API externa \\
S2b & CLI $\rightarrow$ PUB             & Petición-respuesta & Síncrona  & HTTPS/TCP, REST API \\
A1  & PUB $\rightarrow$ Kafka           & Mensajería         & Asíncrona & TCP, Kafka \\
A2  & Kafka $\rightarrow$ MUR           & Mensajería         & Asíncrona & TCP, Kafka \\
S3  & MUR $\rightarrow$ SEG             & RPC                & Síncrona  & HTTP/2, gRPC \\
A3  & Kafka $\rightarrow$ NOT           & Mensajería         & Asíncrona & TCP, Kafka \\
A5  & NOT $\rightarrow$ CLI             & Mensajería push    & Asíncrona & HTTPS/TCP, FCM \\
S10 & CLI $\rightarrow$ MUR             & Petición-respuesta & Síncrona  & HTTPS/TCP, REST API \\
\hline
\end{tabular}
\end{center}

\paragraph{Mecanismos y justificación.}
\begin{itemize}[leftmargin=*]
  \item \textbf{REST sobre HTTPS, síncrono (S1, S2, S2b).} Crear el aviso, subir la imagen
  y vincularla son acciones en las que el usuario necesita una respuesta inmediata
  (\texttt{id\_aviso} o URL) para continuar. HTTPS funciona de forma nativa en Android y en
  el navegador y protege los datos en tránsito.
  \item \textbf{Kafka, asíncrono (A1, A2, A3).} La distribución al muro y las
  notificaciones pueden completarse después de responder al usuario. Kafka desacopla a PUB
  de MUR y de NOT: PUB publica el evento y continúa, y el evento queda guardado aunque un
  consumidor esté caído.
  \item \textbf{gRPC sobre HTTP/2, síncrono (S3).} MUR necesita cada página de seguidores
  para seguir distribuyendo. gRPC trabaja con mensajes binarios compactos y un contrato
  \texttt{.proto} entre servicios internos, por lo que tiene menos latencia que REST con
  JSON.
\end{itemize}

\paragraph{R3: tiempos distintos para avisos normales y alertas.}
Un aviso normal puede tardar hasta 5~s en aparecer en el muro de los seguidores, así que se
acepta consistencia eventual. Una alerta de seguridad debe aparecer en menos de 1~s. Para
que una alerta no quede en cola detrás de los avisos normales, PUB la publica en un tópico
propio (\texttt{alertas.seguridad}) con consumidores dedicados en MUR y en NOT. La alerta
tiene prioridad y capacidad de procesamiento separada del resto.

\paragraph{R4: un negocio con 30.000 seguidores.}
PUB no espera a que el aviso llegue a los 30.000 seguidores, porque la distribución es
asíncrona (A1). MUR pide los seguidores a SEG por páginas (S3) y reparte el trabajo en
lotes entre varios consumidores en paralelo. Además, se limita la carga que puede generar
un solo autor, para que una publicación masiva no ocupe toda la capacidad y el servicio
siga respondiendo al resto de usuarios.

\paragraph{F2: el servicio de muro se cae mientras un negocio publica.}
La distribución no se pierde. MUR confirma un evento en Kafka (\emph{commit} del
\emph{offset}) solo después de procesarlo; si se cae a mitad del trabajo, los eventos no
confirmados siguen en Kafka y, al recuperarse, MUR los vuelve a leer desde el último punto
confirmado. Como un evento puede procesarse dos veces, el consumidor es idempotente:
agregar el mismo aviso al mismo muro no crea un duplicado, por ejemplo gracias a una clave
única (\texttt{id\_usuario}, \texttt{id\_aviso}).

% ------------------------------------------------------------------

\subsection{Publicación y distribución al muro}

\paragraph{Recorrido de un aviso.}
El usuario o el negocio local crea el aviso desde el cliente (CLI). La primera llamada
(S1) envía a PUB el \texttt{id\_usuario}, el texto, la categoría, la ubicación y el tipo
de aviso (normal o alerta); PUB lo registra y responde con el \texttt{id\_aviso}. Si el
aviso lleva imagen, el cliente la sube directamente al servicio externo de almacenamiento
(S2), que devuelve la URL o el identificador del archivo, y después registra esa
referencia en PUB (S2b). Así el archivo no pasa por los servicios internos.

Cuando el aviso está completo, PUB publica el evento \texttt{AvisoPublicado} en Kafka (A1)
y no espera nada más: el usuario ya recibió su confirmación. MUR consume el evento (A2),
consulta a SEG los seguidores del autor mediante gRPC (S3) y agrega el aviso al muro de
cada seguidor. En paralelo, NOT consume el evento \texttt{NuevaPublicación} o
\texttt{AlertaSeguridad} (A3) y envía la notificación push a los seguidores (A5). Cuando un
seguidor abre la aplicación, obtiene de MUR su muro ya armado (S10).

\begin{center}
\small
\setlength{\tabcolsep}{4pt}
\begin{tabular}{|l|l|l|l|l|}
\hline
\textbf{ID} & \textbf{Origen $\rightarrow$ destino} & \textbf{Tipo} & \textbf{Sincronía} & \textbf{Protocolo y middleware} \\
\hline
S1  & CLI $\rightarrow$ PUB             & Petición-respuesta & Síncrona  & HTTPS/TCP, REST API \\
S2  & CLI $\rightarrow$ almac. externo  & Petición-respuesta & Síncrona  & HTTPS/TCP, API externa \\
S2b & CLI $\rightarrow$ PUB             & Petición-respuesta & Síncrona  & HTTPS/TCP, REST API \\
A1  & PUB $\rightarrow$ Kafka           & Mensajería         & Asíncrona & TCP, Kafka \\
A2  & Kafka $\rightarrow$ MUR           & Mensajería         & Asíncrona & TCP, Kafka \\
S3  & MUR $\rightarrow$ SEG             & RPC                & Síncrona  & HTTP/2, gRPC \\
A3  & Kafka $\rightarrow$ NOT           & Mensajería         & Asíncrona & TCP, Kafka \\
A5  & NOT $\rightarrow$ CLI             & Mensajería push    & Asíncrona & HTTPS/TCP, FCM \\
S10 & CLI $\rightarrow$ MUR             & Petición-respuesta & Síncrona  & HTTPS/TCP, REST API \\
\hline
\end{tabular}
\end{center}

\paragraph{Mecanismos y justificación.}
\begin{itemize}[leftmargin=*]
  \item \textbf{REST sobre HTTPS, síncrono (S1, S2, S2b).} Crear el aviso, subir la imagen
  y vincularla son acciones en las que el usuario necesita una respuesta inmediata
  (\texttt{id\_aviso} o URL) para continuar. HTTPS funciona de forma nativa en Android y en
  el navegador y protege los datos en tránsito.
  \item \textbf{Kafka, asíncrono (A1, A2, A3).} La distribución al muro y las
  notificaciones pueden completarse después de responder al usuario. Kafka desacopla a PUB
  de MUR y de NOT: PUB publica el evento y continúa, y el evento queda guardado aunque un
  consumidor esté caído.
  \item \textbf{gRPC sobre HTTP/2, síncrono (S3).} MUR necesita cada página de seguidores
  para seguir distribuyendo. gRPC trabaja con mensajes binarios compactos y un contrato
  \texttt{.proto} entre servicios internos, por lo que tiene menos latencia que REST con
  JSON.
\end{itemize}

\paragraph{R3: tiempos distintos para avisos normales y alertas.}
Un aviso normal puede tardar hasta 5~s en aparecer en el muro de los seguidores, así que se
acepta consistencia eventual. Una alerta de seguridad debe aparecer en menos de 1~s. Para
que una alerta no quede en cola detrás de los avisos normales, PUB la publica en un tópico
propio (\texttt{alertas.seguridad}) con consumidores dedicados en MUR y en NOT. La alerta
tiene prioridad y capacidad de procesamiento separada del resto.

\paragraph{R4: un negocio con 30.000 seguidores.}
PUB no espera a que el aviso llegue a los 30.000 seguidores, porque la distribución es
asíncrona (A1). MUR pide los seguidores a SEG por páginas (S3) y reparte el trabajo en
lotes entre varios consumidores en paralelo. Además, se limita la carga que puede generar
un solo autor, para que una publicación masiva no ocupe toda la capacidad y el servicio
siga respondiendo al resto de usuarios.

\paragraph{F2: el servicio de muro se cae mientras un negocio publica.}
La distribución no se pierde. MUR confirma un evento en Kafka (\emph{commit} del
\emph{offset}) solo después de procesarlo; si se cae a mitad del trabajo, los eventos no
confirmados siguen en Kafka y, al recuperarse, MUR los vuelve a leer desde el último punto
confirmado. Como un evento puede procesarse dos veces, el consumidor es idempotente:
agregar el mismo aviso al mismo muro no crea un duplicado, por ejemplo gracias a una clave
única (\texttt{id\_usuario}, \texttt{id\_aviso}).

% ------------------------------------------------------------------

\subsection{Publicación y distribución al muro}

\paragraph{Recorrido de un aviso.}
El usuario o el negocio local crea el aviso desde el cliente (CLI). La primera llamada
(S1) envía a PUB el \texttt{id\_usuario}, el texto, la categoría, la ubicación y el tipo
de aviso (normal o alerta); PUB lo registra y responde con el \texttt{id\_aviso}. Si el
aviso lleva imagen, el cliente la sube directamente al servicio externo de almacenamiento
(S2), que devuelve la URL o el identificador del archivo, y después registra esa
referencia en PUB (S2b). Así el archivo no pasa por los servicios internos.

Cuando el aviso está completo, PUB publica el evento \texttt{AvisoPublicado} en Kafka (A1)
y no espera nada más: el usuario ya recibió su confirmación. MUR consume el evento (A2),
consulta a SEG los seguidores del autor mediante gRPC (S3) y agrega el aviso al muro de
cada seguidor. En paralelo, NOT consume el evento \texttt{NuevaPublicación} o
\texttt{AlertaSeguridad} (A3) y envía la notificación push a los seguidores (A5). Cuando un
seguidor abre la aplicación, obtiene de MUR su muro ya armado (S10).

\begin{center}
\small
\setlength{\tabcolsep}{4pt}
\begin{tabular}{|l|l|l|l|l|}
\hline
\textbf{ID} & \textbf{Origen $\rightarrow$ destino} & \textbf{Tipo} & \textbf{Sincronía} & \textbf{Protocolo y middleware} \\
\hline
S1  & CLI $\rightarrow$ PUB             & Petición-respuesta & Síncrona  & HTTPS/TCP, REST API \\
S2  & CLI $\rightarrow$ almac. externo  & Petición-respuesta & Síncrona  & HTTPS/TCP, API externa \\
S2b & CLI $\rightarrow$ PUB             & Petición-respuesta & Síncrona  & HTTPS/TCP, REST API \\
A1  & PUB $\rightarrow$ Kafka           & Mensajería         & Asíncrona & TCP, Kafka \\
A2  & Kafka $\rightarrow$ MUR           & Mensajería         & Asíncrona & TCP, Kafka \\
S3  & MUR $\rightarrow$ SEG             & RPC                & Síncrona  & HTTP/2, gRPC \\
A3  & Kafka $\rightarrow$ NOT           & Mensajería         & Asíncrona & TCP, Kafka \\
A5  & NOT $\rightarrow$ CLI             & Mensajería push    & Asíncrona & HTTPS/TCP, FCM \\
S10 & CLI $\rightarrow$ MUR             & Petición-respuesta & Síncrona  & HTTPS/TCP, REST API \\
\hline
\end{tabular}
\end{center}

\paragraph{Mecanismos y justificación.}
\begin{itemize}[leftmargin=*]
  \item \textbf{REST sobre HTTPS, síncrono (S1, S2, S2b).} Crear el aviso, subir la imagen
  y vincularla son acciones en las que el usuario necesita una respuesta inmediata
  (\texttt{id\_aviso} o URL) para continuar. HTTPS funciona de forma nativa en Android y en
  el navegador y protege los datos en tránsito.
  \item \textbf{Kafka, asíncrono (A1, A2, A3).} La distribución al muro y las
  notificaciones pueden completarse después de responder al usuario. Kafka desacopla a PUB
  de MUR y de NOT: PUB publica el evento y continúa, y el evento queda guardado aunque un
  consumidor esté caído.
  \item \textbf{gRPC sobre HTTP/2, síncrono (S3).} MUR necesita cada página de seguidores
  para seguir distribuyendo. gRPC trabaja con mensajes binarios compactos y un contrato
  \texttt{.proto} entre servicios internos, por lo que tiene menos latencia que REST con
  JSON.
\end{itemize}

\paragraph{R3: tiempos distintos para avisos normales y alertas.}
Un aviso normal puede tardar hasta 5~s en aparecer en el muro de los seguidores, así que se
acepta consistencia eventual. Una alerta de seguridad debe aparecer en menos de 1~s. Para
que una alerta no quede en cola detrás de los avisos normales, PUB la publica en un tópico
propio (\texttt{alertas.seguridad}) con consumidores dedicados en MUR y en NOT. La alerta
tiene prioridad y capacidad de procesamiento separada del resto.

\paragraph{R4: un negocio con 30.000 seguidores.}
PUB no espera a que el aviso llegue a los 30.000 seguidores, porque la distribución es
asíncrona (A1). MUR pide los seguidores a SEG por páginas (S3) y reparte el trabajo en
lotes entre varios consumidores en paralelo. Además, se limita la carga que puede generar
un solo autor, para que una publicación masiva no ocupe toda la capacidad y el servicio
siga respondiendo al resto de usuarios.

\paragraph{F2: el servicio de muro se cae mientras un negocio publica.}
La distribución no se pierde. MUR confirma un evento en Kafka (\emph{commit} del
\emph{offset}) solo después de procesarlo; si se cae a mitad del trabajo, los eventos no
confirmados siguen en Kafka y, al recuperarse, MUR los vuelve a leer desde el último punto
confirmado. Como un evento puede procesarse dos veces, el consumidor es idempotente:
agregar el mismo aviso al mismo muro no crea un duplicado, por ejemplo gracias a una clave
única (\texttt{id\_usuario}, \texttt{id\_aviso}).
